% year: תשע"ו
% type: מבחן
% semester: ב
% moed: ב
% source: exams/מבחנים/תשעו/מועד ב תשעו.pdf
% number: 1
% points: 10
% categories: true-false, func-analysis

\begin{question}
\textbf{הוכיחו או הפריכו:} תהי $f$ פונקציה מונוטונית עולה ממש בקטע $[a,b]$ אז $f$ קמורה.
\end{question}

\begin{hint}
חפשו פונקציה עולה ממש שהגרף שלה "מתעקם" לשני הכיוונים, למשל פולינום אי-זוגי.
\end{hint}

\begin{solution}
\textbf{הטענה אינה נכונה.} מונוטוניות אינה קשורה לקמירות.

דוגמה נגדית: $f(x)=x^3$ בקטע $[a,b]=[-1,1]$. מכיוון ש-$f'(x)=3x^2\ge 0$ ומתאפסת רק בנקודה בודדת $x=0$, הפונקציה עולה ממש בקטע (אפשר גם ישירות: אם $x<y$ אז $y^3-x^3=(y-x)(x^2+xy+y^2)>0$, כי $x^2+xy+y^2=(x+\tfrac y2)^2+\tfrac34y^2>0$ כאשר $x\ne y$).

אולם $f''(x)=6x$, כך ש-$f''<0$ ב-$(-1,0)$ ו-$f''>0$ ב-$(0,1)$. כלומר $f$ קעורה ב-$[-1,0]$ וקמורה ב-$[0,1]$, ולכן אינה קמורה (וגם אינה קעורה) בכל הקטע.

בדיקה ישירה לפי ההגדרה: עבור הנקודות $x_1=-1,\ x_2=0$ ונקודת האמצע $-\tfrac12$:
\[ f\!\left(-\tfrac12\right)=-\tfrac18 \;>\; -\tfrac12=\frac{f(-1)+f(0)}{2}, \]
כלומר המיתר בין $(-1,-1)$ ל-$(0,0)$ נמצא \emph{מתחת} לגרף, בסתירה לקמירות ($f(\tfrac{x_1+x_2}{2})\le \tfrac{f(x_1)+f(x_2)}{2}$). ומנגד, עבור $x_1=0,\ x_2=1$: $f(\tfrac12)=\tfrac18<\tfrac12=\tfrac{f(0)+f(1)}2$, כך שהפונקציה גם אינה קעורה. לכן הדוגמה מפריכה את הטענה בכל אחת מהמוסכמות למילה "קמורה".

(דוגמה נוספת: $f(x)=\sqrt{x}$ ב-$[0,1]$ עולה ממש וקעורה ממש.)
\end{solution}
